\section{Actual exceptional normalization retains the true initial homogeneous form}

Fix the actual coordinate-centre blowup chart, with genuine exceptional
variable $x_e$. All normalized members below are the existing internally
constructed quotients of the actual blowup substitution by the stated true
exceptional power; their original definitions are unchanged.

\begin{theorem}[Whole centre precision becomes whole exceptional precision]
\label{hir:actual-normalized-member-exceptional-precision}
\uses{hir:actual-formal-blowup-normalized-member}
\lean{Hironaka.HironakaActualPointExceptionalInitialForm.actualNormalizedFormalMember_sub_mem_exceptional_power}
\leanok
Suppose $F,G\in C^d$ and $F-G\in C^{d+r}$ in the entire actual formal ring,
where $C$ is the coordinate centre ideal. The difference of their actual
$d$-normalized members belongs to the whole exceptional ideal power
$(x_e)^r$, for every $r\in\mathbb N$.
\end{theorem}
\begin{proof}
Let $N_F,N_G$ be the actual $d$-normalized quotients, and internally form
the actual $(d+r)$-normalized quotient $H$ of the whole member $F-G$.
Their genuine product identities give
\[
 x_e^d(N_F-N_G)=\pi^*(F-G)=x_e^{d+r}H
 =x_e^d(x_e^rH).
\]
The actual variable $x_e$ is nonzero in the true formal domain, so genuine
cancellation of $x_e^d$ gives $N_F-N_G=x_e^rH$. This is literal membership
in the whole principal power $(x_e)^r$. The zero precision and zero
normalization degree are included. No equality or quotient certificate is
supplied as an input.
\end{proof}

\begin{lemma}[The entire point-centre remainder has the next true power]
\label{hir:actual-point-homogeneous-remainder-next-power}
\uses{hir:actual-formal-coordinate-power-coefficient-kernel}
\lean{Hironaka.HironakaActualPointExceptionalInitialForm.actual_point_homogeneous_component_and_remainder}
\leanok
For $F\in\mathfrak m^d$, its literal degree-$d$ homogeneous component
$F_d$ belongs to $\mathfrak m^d$ and $F-F_d$ belongs to
$\mathfrak m^{d+1}$, as entire formal ideal memberships.
\end{lemma}
\begin{proof}
The all-variable coordinate weight is total degree. Every coefficient of
$F_d$ outside degree $d$ is zero, proving its first whole membership.
At every exponent of degree below $d+1$, either its degree is $d$ and
the coefficients of $F$ and $F_d$ cancel, or its degree is below $d$
and both are zero by the actual input $F\in\mathfrak m^d$.
The whole coordinate-power coefficient criterion gives the next membership.
\end{proof}

\begin{definition}[The actual finite polynomial of the whole homogeneous component]
\label{hir:actual-point-initial-homogeneous-polynomial}
\lean{Hironaka.HironakaActualPointExceptionalInitialForm.actualPointHomogeneousPolynomial}
\leanok
Set $P_d(F)=\operatorname{trunc}_{<d+1}(F_d)$, using the existing literal
finite total-degree truncation to the original polynomial ring.
\end{definition}

\begin{lemma}[The polynomial inclusion is the entire original homogeneous component]
\label{hir:actual-point-initial-polynomial-whole-inclusion}
\uses{hir:actual-point-initial-homogeneous-polynomial}
\lean{Hironaka.HironakaActualPointExceptionalInitialForm.actualPointHomogeneousPolynomial_coe}
\leanok
The genuine polynomial inclusion of $P_d(F)$ equals the entire series $F_d$.
\end{lemma}
\begin{proof}
At each exponent of degree below $d+1$, the actual truncation retains its
coefficient. At any remaining exponent the degree differs from $d$, so
the actual homogeneous-component coefficient is zero. Equality of every
coefficient proves the entire polynomial-inclusion identity.
\end{proof}

\begin{theorem}[The actual exceptional restriction depends only on the true homogeneous initial polynomial]
\label{hir:actual-point-exceptional-initial-polynomial-restriction}
\uses{hir:actual-normalized-member-exceptional-precision,hir:actual-point-homogeneous-remainder-next-power,hir:actual-point-initial-polynomial-whole-inclusion}
\lean{Hironaka.HironakaActualPointExceptionalInitialForm.actual_point_normalized_member_congr_homogeneous_polynomial}
\leanok
For every actual $F\in\mathfrak m^d$, the actual $d$-normalized member
of $F$ and the actual $d$-normalized member of the entire polynomial
inclusion of $P_d(F)$ differ by a member of $(x_e)$. Thus they have the
same whole class on the exceptional hyperplane in this true point chart.
\end{theorem}
\begin{proof}
Use the preceding entire polynomial-inclusion identity to supply its
actual point-power membership. The genuine remainder theorem makes the
difference of the two original whole numerators a member of
$\mathfrak m^{d+1}$. Apply the general exceptional-precision theorem
with $r=1$. The resulting principal first power is exactly $(x_e)$.
This compares the actual normalized quotients, including every coefficient
in the entire exceptional restriction, without inferring an infinite
equality from a fixed finite jet.
\end{proof}

This is the initial-form interface for the point-centre Bennett argument.
The dehomogenized polynomial comparison and the resulting all-order
Hilbert--Samuel inequality are further obligations, not consequences
claimed here. General positive-dimensional permissible centres remain open.
